
\begin{equation}\tag{6.1.1}
 M^{\underline{n}},\qquad \bigotimes^{\underline{n}} M,\qquad \bigwedge^{\underline{n}} M,\quad\text{etc.}
\end{equation}
Les
foncteurs \((\ )^{\underline{n}}\), \(\bigotimes\limits^{\underline{n}}\) et \(\bigwedge\limits^{\underline{n}}\) transforment modules libres (resp. localement libres) en modules libres (resp. localement libres).

De façon générale, considérons des sections \(\underline{n}_1,\ldots,n_r\) de \(\mathbb Z_X\) et un multifoncteur (covariant, pour fixer les idées)
\[
 T:\underline{\Mod}^{n_1}(\underline{O})\times\cdots\times
 \underline{\Mod}^{n_r}(\underline{O})\longrightarrow\underline{\Mod}(\underline{O})
\]
à \(r\) variables, où \(\Mod_{\underline{O}}^{\underline{n}}(X)\) désigne la catégorie fibrée sur \(X\) des \(\underline{O}\)-modules localement libres de rang \(\underline{n}\). Soit
\[
 M=T(\underline{O}^{\underline{n}_1},\ldots,\underline{O}^{\underline{n}_r})
\]
et
\[
 G=\operatorname{G\ell}(\underline{n}_1)\times\cdots\times\operatorname{G\ell}(\underline{n}_r),
 \quad\text{où }\operatorname{G\ell}(\underline{n}_i)=\underline{\Aut}_{\underline{O}}(\underline{O}^{n_i}).
\]
Il s'ensuit immédiatement que \(T\) définit une représentation
\begin{equation}\tag{6.1.2}
 \alpha_T:G\longrightarrow\underline{\Aut}_{\underline{O}}(M).
\end{equation}
Il est clair que si \(M\) est lui-même localement libre de rang \(\underline{m}\), pour tout foncteur composé,
\[
 T'T,\ \text{où } T':\underline{\Mod}^{\underline{m}}(\underline{O})\longrightarrow\underline{\Mod}(\underline{O}),
\]
on a
\begin{equation}\tag{6.1.3}
 \alpha_{T'T}=\alpha_{T'}\,\alpha_T.
\end{equation}

Réciproquement, étant donné la représentation \(\alpha=\alpha_T:G\to\underline{\Aut}_{\underline{O}}(M)\), on récupère \(T(\underline{E}_1,\ldots,\underline{E}_r)\) (pour toute suite
